% 2de graad, reeks Vectoren (Vectoren, Analytische meetkunde).
% Signac: zeilboten; de wind duwt met een richting en een grootte.
% Motief: twee vectoren en hun som, volgens de parallellogramregel.
\def\omslagbestand{omslag-saint-tropez.jpg}
\def\omslagbron{Paul Signac, \textit{De haven van Saint-Tropez}, 1901-1902 (publiek domein)}
\def\omslagrand{10.6 + 0.3*sin(\x*360/7)}
\def\omslagmotief{%
  \randlijn
  \draw[white,line width=.7pt,dashed] (17.2,5.8) -- (18.1,8.1) -- (15.5,7.2);
  \draw[white,line width=1.4pt,-{Stealth[length=3mm]}] (14.6,4.9) -- (17.2,5.8);
  \draw[white,line width=1.4pt,-{Stealth[length=3mm]}] (14.6,4.9) -- (15.5,7.2);
  \draw[fel,line width=2pt,-{Stealth[length=3.4mm]}] (14.6,4.9) -- (18.1,8.1);
  \node[text=white,anchor=north] at (16.1,5.25) {$\vec u$};
  \node[text=white,anchor=east] at (14.95,6.2) {$\vec v$};
  \node[text=fel,anchor=west] at (18.25,8.25) {$\vec u+\vec v$};
}
